\documentclass[10pt,a4paper]{beamer}
\usetheme{Singapore}
\setbeamercovered{transparent}
\beamertemplatenavigationsymbolsempty % Comment this out if you want a navigation bar

\usepackage[utf8]{inputenc}
\usepackage[english]{babel}

%math packages
\usepackage{amsmath}
\usepackage{amsfonts}
\usepackage{amssymb}

%package for rule inferences
\usepackage{bussproofs}

%package for handling verbatim
\usepackage{fancyvrb}

% packages for drawing
\usepackage{pgf}
\usepackage{tikz}
\usetikzlibrary{arrows,matrix}

% package for strikethrough text
\usepackage{ulem}

\logo{
	\includegraphics[height=1cm]{coq_logo.png}
}

% uncover everything in a step-wise fashion
%\beamerdefaultoverlayspecification{<+->}

\usepackage{color}
\usepackage{xcolor}

\definecolor{coqGreen}{rgb}{0,0.5,0}

% Environment for user input / source code
\DefineVerbatimEnvironment
{user}{Verbatim}
{formatcom=\color{coqGreen}}

% Environment for Coq output
\DefineVerbatimEnvironment
{coq}{Verbatim}
{formatcom=\color{blue}}

% Customizable emphasis
\providecommand{\highlight}[1]{\textcolor{red}{#1}}
\providecommand{\Highlight}[1]{\textcolor{red}{\textbf{#1}}}
\title{Coq Proof Assistant: Propositions and Proofs}
\author{Mirco Kocher}
\institute{
	Logic and Theory Group\\
	Institute of Computer Science and Applied Mathematics\\
	Universität Bern
}
\date{2012}
\begin{document}

{ \logo{} % Suppress logo for title page
\begin{frame} % Title page
	\begin{center}
		\includegraphics{coq_logo.png}
	\end{center}	
	\titlepage	
\end{frame}
}

\begin{frame}{Overview}
\tableofcontents[pausesections, pausesubsections]
\end{frame}

\section{Minimal Propositional Logic}

\subsection{Basics}
\begin{frame}{Truth table}
	$(P \rightarrow Q)$
	\begin{itemize}
		\item Classical logic
		\pause
		\item Assign to every variable a denotation true or false
		\pause		
		\item Formula is valid iff true in all cases
		\pause
		\item Question "is the proposition P true?"
	\end{itemize}
\end{frame}
\begin{frame}{Coq system}
	$(P \rightarrow Q)$
	\begin{itemize}
		\item Intuitionistic logic
		\pause
		\item Obtain a proof of Q from a proof of P
		\pause		
		\item Arbitrary proof of P constructs a proof of Q
		\pause
		\item Question "what are the proof of P (if any)?"
	\end{itemize}
\end{frame}

\subsection{Definition}
\begin{frame}[fragile]{Hypothesis}
	\begin{user}
	Hypothesis h:P
	\end{user}
	\begin{itemize}
		\item Local declaration
		\pause
		\item $h$ is the name of the hypothesis
		\pause
		\item P is its statement
		\pause
		\item Synonymous to
		\begin{user}
		Variable h:P
		\end{user}
		\pause
		\item Use
		\begin{user}
		Hypotheses
		\end{user}
		or
		\begin{user}
		Variables
		\end{user}
		to declare several at a time
	\end{itemize}
\end{frame}
\begin{frame}[fragile]{Section}
	The section contains all Hypoteses / Variables from the Context\\
	\pause
	\medskip
	Start section sec1 with
	\begin{user}
	Section sec1
	\end{user}
	\pause
	\medskip
	End section sec1 with
	\begin{user}
	End sec1
	\end{user}
\end{frame}
\begin{frame}[fragile]{Axiom}
	\begin{user}
	Axiom x:P
	\end{user}
	\begin{itemize}
		\item Global declaration
		\pause
		\item Synonymous to
		\begin{user}
		Parameter x:P
		\end{user}
	\end{itemize}
	\pause
	\medskip
	Environment contains axioms\\
	\pause
	\medskip
	Context contains hypotheses\\
	\pause
	\medskip
	$E, \Gamma \vdash \pi : P$
\end{frame}
\begin{frame}{Goals and Tactics}
	Goals: what needs to be proven\\
	\medskip
	\pause
	Goal: $E, \Gamma \vdash P$\\
	\pause
	Construct a proof of P. Should be a well-formed term $t$ in the environment $E$ and context $\Gamma$\\
	\pause
	Term $t$ is called a $solution$\\
	\pause
	\bigskip
	Tactics: commands to decompose this goal into simpler goals\\
	\pause
	\medskip
	$g$ ist input goal and $g_1$, $g_2$, ..., $g_k$ are output goals\\
	\pause
	Possible to construct a solution of $g$ from the solutions of goals $g_i$	
\end{frame}

\section{Example}

\subsection{Definition}
\begin{frame}[fragile]{intros}
	Goal: $(P \rightarrow Q) \rightarrow (Q \rightarrow R) \rightarrow (P \rightarrow R)$ \\
	\medskip
	\pause
	\begin{user}
	intros H H' p
	\end{user}
	\pause
	\begin{itemize}
		\item Transform the task of construction a proof into proving R with those hypotheses added
		\pause
		\item $H:P \rightarrow Q$
		\pause
		\item $H':Q \rightarrow R$ and
		\pause
		\item $p:P$
		\pause
		\item New subgoal: $R$
	\end{itemize}
	\pause
	Simplifies the statement to prove and increases the resources available
\end{frame}
\begin{frame}[fragile]{apply}
	Subgoal: $R$\\
	\medskip
	Hypothesis $H': Q \rightarrow R$ and $H:P \rightarrow Q$
	\medskip
	\pause
	\begin{user}
	apply H'
	\end{user}
	\begin{itemize}
		\item Use the hypothesis $H'$ to advance our proof
		\pause
		\item Argument has to be a premise and a conclusion
		\pause
		\item Creates new goal for the premise
		\pause
		\item New subgoal: $Q$
	\end{itemize}
	\pause
	Applying hypothesis H gives the new subgoal $P$
\end{frame}
\begin{frame}[fragile]{assumption}
	Subgoal: $P$\\
	\medskip
	Hypothesis $p: P$
	\medskip
	\pause
	\begin{user}
	assumption
	\end{user}
	\begin{itemize}
		\item Statement to proof is exactly statement of hypothesis $p$
		\pause
		\item Succeeds without generating any new goal
	\end{itemize}
	\pause
	\medskip
	\begin{coq}
	No more subgoals
	\end{coq}
	\pause
	Proof complete
\end{frame}
\begin{frame}[fragile]{Finish}
	\begin{user}
	Qed
	\end{user}
	\pause
	\begin{itemize}
		\item Saves the theorem's name, statement and proof term
		\pause
		\item Displays the sequence of tactics.
	\end{itemize}
	\begin{coq}
	intros H H' p.
	apply H'.
	apply H.
	assumption.
	\end{coq}
	\pause
	\medskip
	\begin{user}
	Print theorem-name
	\end{user}
	\pause
	\begin{itemize}
		\item Shows the proof like any $Gallina$ definition
	\end{itemize}
	\begin{coq}
	theorem-name = fun (H:P -> Q)(H':Q -> R)(p:P) => H' (H p)
		 : (P -> Q) -> (Q -> R) -> P -> R
	\end{coq}
\end{frame}

\subsection{Demo}
\begin{frame}{Transitivity}
	\begin{center}
		$(P \rightarrow Q) \rightarrow (Q \rightarrow R) \rightarrow (P \rightarrow R)$
	\end{center}
\end{frame}

\section{Structure}

\subsection{Tactics explenation}
\begin{frame}[fragile]{One Shot}
	Not all the details for the proof is needed\\
	\medskip
	\pause
	A few automatic tactics are able to solve this goal
	\pause
	\begin{user}
	Theorem transitivity : (P->Q) -> (Q->R) -> P -> R.
	Proof.
	\end{user}
	\pause
	\begin{user}
		auto.
	Qed.
	\end{user}
	\pause
	Good usage requires the command Proof after Theorem or Lemma\\
	\smallskip
	Makes the proof documents more readable
\end{frame}
\begin{frame}[fragile]{exact}
	The exact tactic takes any term (with right type) as argument\\
	\pause
	\smallskip
	The whole proof could be given as an argument\\
	\pause
	\medskip
	\begin{user}
	Theorem delta : (P->P->Q) -> P -> Q.
	Proof.
	\end{user}
	\pause
	\begin{user}
	exact (fun (H:P->P->Q)(p:P) => H p p).
	Qed.
	\end{user}
	\pause
	\medskip
	Or even shorter:
	\begin{user}
	Theorem delta : (P->P->Q) -> P -> Q.
	Proof (fun (H:P->P->Q)(p:P) => H p p).
	\end{user}
\end{frame}
\begin{frame}[fragile]{Modus Ponens}
	apply tactic uses the Modus Ponens\\
	\pause
	\medskip
	\begin{prooftree}
	\AxiomC{$E,\Gamma \vdash t : P \rightarrow Q$}
	\AxiomC{$E,\Gamma \vdash t' : P$}
	\BinaryInfC{$E, \Gamma \vdash t t' : Q$}\textbf{App}
	\end{prooftree}
	\pause
	Term t : $P_1 \rightarrow P_2 \rightarrow ... \rightarrow P_n \rightarrow Q$ \\
	Goal : $P_k \rightarrow P_{k+1} \rightarrow ... \rightarrow P_n \rightarrow Q$
	\pause
	\medskip
	\begin{user}
	apply t
	\end{user}
	generates k-1 goals with statements $P_1, ..., P_{k-1}$\\
	\pause
	\smallskip
	if new goals have solution $t_1, t_2, ..., t_{k-1}$\\
	solution is $t$ $t_1$ $t_2$ ... $t_{k-1}$ for initial goal
\end{frame}
\begin{frame}[fragile]{intros}
	\begin{user}
	intros v_1, v_2, ..., v_n
	\end{user}
	is the same as
	\begin{user}
	intro v_1, intro v_2, ..., intro v_n
	\end{user}
	\pause
	intro $v_1$ takes the first implication as a hypothesis called $v_1$\\
	\bigskip
	\pause
	Proof Theorem K : $P \rightarrow Q \rightarrow P$\\
	\pause
	Goal is  $P \rightarrow Q \rightarrow P$\\
	\pause
	\begin{user}
	intro p.
	\end{user}
	Hypothesis $p:P$ and new Goal $Q \rightarrow P$
\end{frame}
\begin{frame}[fragile]{Handling}
	\begin{user}
		Show i
	\end{user}
	Display goal $i$ with complete context\\
	Coq displays the goals after each proof step
	\pause
	\medskip
	\begin{user}
		Undo n
	\end{user}
	Go back $n$ steps and try an alternative if goal can not be solved
	\pause
	\medskip
	\begin{user}
		Restart
	\end{user}
	Go back to the beginning of the proof
	\pause
	\medskip
	\begin{user}
		Abort
	\end{user}
	Abandon the proof
\end{frame}

\subsection{Composing tactics}
\begin{frame}[fragile]{Simple composing}
	Combine tactics without stopping at intermediary subgoals\\
	\pause
	\medskip
	Goal: $P \rightarrow Q \rightarrow (P \rightarrow Q \rightarrow R) \rightarrow R$
	\pause
	\begin{user}
	intros p q H; apply H; assumption.
	\end{user}
	\pause
	Like chess: forsee results of tactics\\
	\pause
	\medskip
	If any tatcic fails, then the whole combination fails
\end{frame}
\begin{frame}[fragile]{General composing}
	Tactics can generate multiple subgoals \\
	\pause
	\medskip
	Goal: $(P \rightarrow Q \rightarrow R) \rightarrow (P \rightarrow Q) \rightarrow (P \rightarrow R)$
	\pause
	\begin{user}
	intros H H' p; apply H; [assumption | apply H'; assumption].
	\end{user}
	\pause
	Two subgoals $P$ and $Q$\\
	\pause
	\medskip
	First is solved with assumption\\
	\pause
	Other one first has to apply H' and then use assumption
\end{frame}
\begin{frame}[fragile]{More composing}
	If a tactic fails automatically use another tactic
	\pause
	\begin{user}
	intros H p; apply H; (assumption || intro H').
	\end{user}
	\pause
	If assumption fails, then intro H' is used\\
	\pause
	\bigskip
	A tactic can be left unchanged to finish every subgoal in one go\\
	\pause
	\smallskip
	Goal: $(P \rightarrow Q) \rightarrow (P \rightarrow R) \rightarrow (P \rightarrow Q \rightarrow R \rightarrow T) \rightarrow P \rightarrow T$
	\pause
	\begin{user}
		intros H H0 H1 p.
		apply H1; [idtac | apply H | apply H0]; assumption
	\end{user}
\end{frame}
\begin{frame}[fragile]{Fail}
	Tactic that always fails \\
	\pause
	\medskip
	Goal: $(P \rightarrow Q) \rightarrow (P \rightarrow Q)$
	\pause
	\begin{user}
	intro X; apply X; fail.
	\end{user}
	\pause
	This combination succeeds; there are no more subgoals after "apply X"\\
	\pause
	\medskip
	Goal: $((P \rightarrow P) \rightarrow (Q \rightarrow Q) \rightarrow R) \rightarrow R$
	\begin{user}
	intro X; apply X; fail.
	\end{user}
	\pause
	This combination fails; there are subgoals left after "apply X"
\end{frame}
\begin{frame}[fragile]{Try}
	Combination of tactics that never fail \\
	\pause
	\medskip
	Goal: $(P \rightarrow Q \rightarrow R \rightarrow T) \rightarrow (P \rightarrow Q) \rightarrow (P \rightarrow R \rightarrow T)$
	\pause
	\begin{user}
	intros H H' p r.
	apply H; try assumption; apply H'; assumption.
	\end{user}
	\pause
	"try tac" behaves like "tac $\|$ idtac"\\
	\pause
	\medskip
	tac is either applied or the subgoal is left unchanged
\end{frame}
\begin{frame}[fragile]{Unprovalbe Propositions}
	There are goals with no solution at all\\
	\pause
	\medskip
	Even though they are valid in classical logic\\
	\pause
	\medskip
	Peirce's formula: $(((P \rightarrow Q) \rightarrow P) \rightarrow P)$\\
	\pause
	\medskip
	Truth table shows it is a valid formula
\end{frame}

\end{document}